\documentclass[10pt,a4paper]{amsart}
\usepackage{amsmath,amssymb,amsthm,graphicx}
% For commutative diagrams we use Xy-pic
\usepackage[all]{xy}

% We use 2cell for 2-commutative diagrams.
\xyoption{2cell}
% Theorem environments.
%
\theoremstyle{plain}
\newtheorem{theorem}[subsection]{Theorem}
\newtheorem{proposition}[subsection]{Proposition}
\newtheorem{lemma}[subsection]{Lemma}

\theoremstyle{definition}
\newtheorem{definition}[subsection]{Definition}
\newtheorem{example}[subsection]{Example}
\newtheorem{exercise}[subsection]{Exercise}
\newtheorem{situation}[subsection]{Situation}

\theoremstyle{remark}
\newtheorem{remark}[subsection]{Remark}
\newtheorem{remarks}[subsection]{Remarks}

\numberwithin{equation}{subsection}

% Macros
%
\def\lim{\mathop{\mathrm{lim}}\nolimits}
\def\colim{\mathop{\mathrm{colim}}\nolimits}
\def\Spec{\mathop{\mathrm{Spec}}}
\def\Hom{\mathop{\mathrm{Hom}}\nolimits}
\def\Ext{\mathop{\mathrm{Ext}}\nolimits}
\def\SheafHom{\mathop{\mathcal{H}\!\mathit{om}}\nolimits}
\def\SheafExt{\mathop{\mathcal{E}\!\mathit{xt}}\nolimits}
\def\Sch{\mathit{Sch}}
\def\Mor{\mathop{\mathrm{Mor}}\nolimits}
\def\Ob{\mathop{\mathrm{Ob}}\nolimits}
\def\Sh{\mathop{\mathit{Sh}}\nolimits}
\def\NL{\mathop{N\!L}\nolimits}
\def\CH{\mathop{\mathrm{CH}}\nolimits}
\def\proetale{{pro\text{-}\acute{e}tale}}
\def\etale{{\acute{e}tale}}
\def\QCoh{\mathit{QCoh}}
\def\Ker{\mathop{\mathrm{Ker}}}
\def\Im{\mathop{\mathrm{Im}}}
\def\Coker{\mathop{\mathrm{Coker}}}
\def\Coim{\mathop{\mathrm{Coim}}}

% Boxtimes
%
\DeclareMathSymbol{\boxtimes}{\mathbin}{AMSa}{"02}

%
% Macros for moduli stacks/spaces
%
\def\QCohstack{\mathcal{QC}\!\mathit{oh}}
\def\Cohstack{\mathcal{C}\!\mathit{oh}}
\def\Spacesstack{\mathcal{S}\!\mathit{paces}}
\def\Quotfunctor{\mathrm{Quot}}
\def\Hilbfunctor{\mathrm{Hilb}}
\def\Curvesstack{\mathcal{C}\!\mathit{urves}}
\def\Polarizedstack{\mathcal{P}\!\mathit{olarized}}
\def\Complexesstack{\mathcal{C}\!\mathit{omplexes}}
% \Pic is the operator that assigns to X its picard group, usage \Pic(X)
% \Picardstack_{X/B} denotes the Picard stack of X over B
% \Picardfunctor_{X/B} denotes the Picard functor of X over B
\def\Pic{\mathop{\mathrm{Pic}}\nolimits}
\def\Picardstack{\mathcal{P}\!\mathit{ic}}
\def\Picardfunctor{\mathrm{Pic}}
\def\Deformationcategory{\mathcal{D}\!\mathit{ef}}
\newenvironment{reference}{}{}
\newenvironment{slogan}{}{}
\newenvironment{history}{}{}
\usepackage{cleveref}
\crefname{originalsection}{original section}{original sections}
\Crefname{originalsection}{Original section}{Original sections}
\makeatletter
\newlabel{groupoids-quotients-section-introduction}{{048B}{source}{}{sourcepointer-1}{}}
\newlabel{groupoids-quotients-section-introduction@cref}{{[originalsection][0][]048B}{[1][0][]source}{}{}{}}
\makeatother
\crefname{originaltheorem}{original theorem}{original theorems}
\Crefname{originaltheorem}{Original theorem}{Original theorems}
\makeatletter
\newlabel{criteria-theorem-flat-groupoid-gives-algebraic-stack}{{06FI}{source}{}{sourcepointer-2}{}}
\newlabel{criteria-theorem-flat-groupoid-gives-algebraic-stack@cref}{{[originaltheorem][0][]06FI}{[1][0][]source}{}{}{}}
\makeatother
\crefname{originallemma}{original lemma}{original lemmas}
\Crefname{originallemma}{Original lemma}{Original lemmas}
\makeatletter
\newlabel{spaces-groupoids-lemma-quotient-stack-2-cartesian}{{04M9}{source}{}{sourcepointer-3}{}}
\newlabel{spaces-groupoids-lemma-quotient-stack-2-cartesian@cref}{{[originallemma][0][]04M9}{[1][0][]source}{}{}{}}
\makeatother
\crefname{originallemma}{original lemma}{original lemmas}
\Crefname{originallemma}{Original lemma}{Original lemmas}
\makeatletter
\newlabel{spaces-groupoids-lemma-quotient-stack-2-coequalizer}{{044U}{source}{}{sourcepointer-4}{}}
\newlabel{spaces-groupoids-lemma-quotient-stack-2-coequalizer@cref}{{[originallemma][0][]044U}{[1][0][]source}{}{}{}}
\makeatother
\crefname{originallemma}{original lemma}{original lemmas}
\Crefname{originallemma}{Original lemma}{Original lemmas}
\makeatletter
\newlabel{spaces-groupoids-lemma-quotient-stack-objects}{{044X}{source}{}{sourcepointer-5}{}}
\newlabel{spaces-groupoids-lemma-quotient-stack-objects@cref}{{[originallemma][0][]044X}{[1][0][]source}{}{}{}}
\makeatother
\crefname{originaldefinition}{original definition}{original definitions}
\Crefname{originaldefinition}{Original definition}{Original definitions}
\makeatletter
\newlabel{groupoids-quotients-definition-categorical}{{048J}{source}{}{sourcepointer-6}{}}
\newlabel{groupoids-quotients-definition-categorical@cref}{{[originaldefinition][0][]048J}{[1][0][]source}{}{}{}}
\makeatother
\crefname{originallemma}{original lemma}{original lemmas}
\Crefname{originallemma}{Original lemma}{Original lemmas}
\makeatletter
\newlabel{criteria-lemma-flat-quotient-flat-presentation}{{06FH}{source}{}{sourcepointer-7}{}}
\newlabel{criteria-lemma-flat-quotient-flat-presentation@cref}{{[originallemma][0][]06FH}{[1][0][]source}{}{}{}}
\makeatother
\crefname{originalsection}{original section}{original sections}
\Crefname{originalsection}{Original section}{Original sections}
\makeatletter
\newlabel{algebraic-remark-flat-fp-presentation}{{04WY}{source}{}{sourcepointer-8}{}}
\newlabel{algebraic-remark-flat-fp-presentation@cref}{{[originalsection][0][]04WY}{[1][0][]source}{}{}{}}
\makeatother
\crefname{originalsection}{original section}{original sections}
\Crefname{originalsection}{Original section}{Original sections}
\makeatletter
\newlabel{spaces-section-presentations}{{0261}{source}{}{sourcepointer-9}{}}
\newlabel{spaces-section-presentations@cref}{{[originalsection][0][]0261}{[1][0][]source}{}{}{}}
\makeatother
\crefname{originalsection}{original section}{original sections}
\Crefname{originalsection}{Original section}{Original sections}
\makeatletter
\newlabel{guide-subsection-coarse-moduli-spaces}{{04UX}{source}{}{sourcepointer-10}{}}
\newlabel{guide-subsection-coarse-moduli-spaces@cref}{{[originalsection][0][]04UX}{[1][0][]source}{}{}{}}
\makeatother
\crefname{originallemma}{original lemma}{original lemmas}
\Crefname{originallemma}{Original lemma}{Original lemmas}
\makeatletter
\newlabel{algebraic-lemma-map-space-into-stack}{{04T4}{source}{}{sourcepointer-11}{}}
\newlabel{algebraic-lemma-map-space-into-stack@cref}{{[originallemma][0][]04T4}{[1][0][]source}{}{}{}}
\makeatother
\crefname{originallemma}{original lemma}{original lemmas}
\Crefname{originallemma}{Original lemma}{Original lemmas}
\makeatletter
\newlabel{morphisms-lemma-finite-flat}{{02KB}{source}{}{sourcepointer-12}{}}
\newlabel{morphisms-lemma-finite-flat@cref}{{[originallemma][0][]02KB}{[1][0][]source}{}{}{}}
\makeatother
\crefname{originallemma}{original lemma}{original lemmas}
\Crefname{originallemma}{Original lemma}{Original lemmas}
\makeatletter
\newlabel{stacks-properties-lemma-check-property-covering}{{04XD}{source}{}{sourcepointer-13}{}}
\newlabel{stacks-properties-lemma-check-property-covering@cref}{{[originallemma][0][]04XD}{[1][0][]source}{}{}{}}
\makeatother
\crefname{originallemma}{original lemma}{original lemmas}
\Crefname{originallemma}{Original lemma}{Original lemmas}
\makeatletter
\newlabel{stacks-properties-lemma-quasi-compact-stack}{{04YC}{source}{}{sourcepointer-14}{}}
\newlabel{stacks-properties-lemma-quasi-compact-stack@cref}{{[originallemma][0][]04YC}{[1][0][]source}{}{}{}}
\makeatother
\crefname{originallemma}{original lemma}{original lemmas}
\Crefname{originallemma}{Original lemma}{Original lemmas}
\makeatletter
\newlabel{stacks-morphisms-lemma-properties-diagonal-from-presentation}{{0DTY}{source}{}{sourcepointer-15}{}}
\newlabel{stacks-morphisms-lemma-properties-diagonal-from-presentation@cref}{{[originallemma][0][]0DTY}{[1][0][]source}{}{}{}}
\makeatother
\crefname{originallemma}{original lemma}{original lemmas}
\Crefname{originallemma}{Original lemma}{Original lemmas}
\makeatletter
\newlabel{spaces-limits-lemma-affine}{{07VS}{source}{}{sourcepointer-16}{}}
\newlabel{spaces-limits-lemma-affine@cref}{{[originallemma][0][]07VS}{[1][0][]source}{}{}{}}
\makeatother
\crefname{originalproposition}{original proposition}{original propositions}
\Crefname{originalproposition}{Original proposition}{Original propositions}
\makeatletter
\newlabel{groupoids-proposition-finite-flat-equivalence}{{03BM}{source}{}{sourcepointer-17}{}}
\newlabel{groupoids-proposition-finite-flat-equivalence@cref}{{[originalproposition][0][]03BM}{[1][0][]source}{}{}{}}
\makeatother
\crefname{originallemma}{original lemma}{original lemmas}
\Crefname{originallemma}{Original lemma}{Original lemmas}
\makeatletter
\newlabel{stacks-morphisms-lemma-base-change-affine}{{0CHR}{source}{}{sourcepointer-18}{}}
\newlabel{stacks-morphisms-lemma-base-change-affine@cref}{{[originallemma][0][]0CHR}{[1][0][]source}{}{}{}}
\makeatother
\crefname{originalsection}{original section}{original sections}
\Crefname{originalsection}{Original section}{Original sections}
\makeatletter
\newlabel{groupoids-section-finite-flat}{{03BE}{source}{}{sourcepointer-19}{}}
\newlabel{groupoids-section-finite-flat@cref}{{[originalsection][0][]03BE}{[1][0][]source}{}{}{}}
\makeatother
\crefname{originallemma}{original lemma}{original lemmas}
\Crefname{originallemma}{Original lemma}{Original lemmas}
\makeatletter
\newlabel{schemes-lemma-morphism-into-affine}{{01I1}{source}{}{sourcepointer-20}{}}
\newlabel{schemes-lemma-morphism-into-affine@cref}{{[originallemma][0][]01I1}{[1][0][]source}{}{}{}}
\makeatother
\crefname{originallemma}{original lemma}{original lemmas}
\Crefname{originallemma}{Original lemma}{Original lemmas}
\makeatletter
\newlabel{stacks-morphisms-lemma-compose-after-separated}{{050M}{source}{}{sourcepointer-21}{}}
\newlabel{stacks-morphisms-lemma-compose-after-separated@cref}{{[originallemma][0][]050M}{[1][0][]source}{}{}{}}
\makeatother
\crefname{originallemma}{original lemma}{original lemmas}
\Crefname{originallemma}{Original lemma}{Original lemmas}
\makeatletter
\newlabel{stacks-morphisms-lemma-surjection-from-quasi-compact}{{050X}{source}{}{sourcepointer-22}{}}
\newlabel{stacks-morphisms-lemma-surjection-from-quasi-compact@cref}{{[originallemma][0][]050X}{[1][0][]source}{}{}{}}
\makeatother
\crefname{originallemma}{original lemma}{original lemmas}
\Crefname{originallemma}{Original lemma}{Original lemmas}
\makeatletter
\newlabel{groupoids-lemma-integral-over-invariants}{{03BJ}{source}{}{sourcepointer-23}{}}
\newlabel{groupoids-lemma-integral-over-invariants@cref}{{[originallemma][0][]03BJ}{[1][0][]source}{}{}{}}
\makeatother
\crefname{originallemma}{original lemma}{original lemmas}
\Crefname{originallemma}{Original lemma}{Original lemmas}
\makeatletter
\newlabel{morphisms-lemma-integral-universally-closed}{{01WM}{source}{}{sourcepointer-24}{}}
\newlabel{morphisms-lemma-integral-universally-closed@cref}{{[originallemma][0][]01WM}{[1][0][]source}{}{}{}}
\makeatother
\crefname{originallemma}{original lemma}{original lemmas}
\Crefname{originallemma}{Original lemma}{Original lemmas}
\makeatletter
\newlabel{stacks-morphisms-lemma-image-proper-is-proper}{{0CQK}{source}{}{sourcepointer-25}{}}
\newlabel{stacks-morphisms-lemma-image-proper-is-proper@cref}{{[originallemma][0][]0CQK}{[1][0][]source}{}{}{}}
\makeatother
\crefname{originallemma}{original lemma}{original lemmas}
\Crefname{originallemma}{Original lemma}{Original lemmas}
\makeatletter
\newlabel{stacks-morphisms-lemma-points-presentation}{{0DTZ}{source}{}{sourcepointer-26}{}}
\newlabel{stacks-morphisms-lemma-points-presentation@cref}{{[originallemma][0][]0DTZ}{[1][0][]source}{}{}{}}
\makeatother
\crefname{originallemma}{original lemma}{original lemmas}
\Crefname{originallemma}{Original lemma}{Original lemmas}
\makeatletter
\newlabel{groupoids-lemma-points}{{03BL}{source}{}{sourcepointer-27}{}}
\newlabel{groupoids-lemma-points@cref}{{[originallemma][0][]03BL}{[1][0][]source}{}{}{}}
\makeatother
\crefname{originallemma}{original lemma}{original lemmas}
\Crefname{originallemma}{Original lemma}{Original lemmas}
\makeatletter
\newlabel{groupoids-quotients-lemma-base-change-quotient-stack}{{0DTG}{source}{}{sourcepointer-28}{}}
\newlabel{groupoids-quotients-lemma-base-change-quotient-stack@cref}{{[originallemma][0][]0DTG}{[1][0][]source}{}{}{}}
\makeatother
\crefname{originallemma}{original lemma}{original lemmas}
\Crefname{originallemma}{Original lemma}{Original lemmas}
\makeatletter
\newlabel{groupoids-lemma-invariants-base-change}{{03BK}{source}{}{sourcepointer-29}{}}
\newlabel{groupoids-lemma-invariants-base-change@cref}{{[originallemma][0][]03BK}{[1][0][]source}{}{}{}}
\makeatother
\crefname{originallemma}{original lemma}{original lemmas}
\Crefname{originallemma}{Original lemma}{Original lemmas}
\makeatletter
\newlabel{algebra-lemma-universally-bijective}{{09EF}{source}{}{sourcepointer-30}{}}
\newlabel{algebra-lemma-universally-bijective@cref}{{[originallemma][0][]09EF}{[1][0][]source}{}{}{}}
\makeatother
\crefname{originallemma}{original lemma}{original lemmas}
\Crefname{originallemma}{Original lemma}{Original lemmas}
\makeatletter
\newlabel{stacks-morphisms-lemma-universally-injective-local}{{0DTN}{source}{}{sourcepointer-31}{}}
\newlabel{stacks-morphisms-lemma-universally-injective-local@cref}{{[originallemma][0][]0DTN}{[1][0][]source}{}{}{}}
\makeatother
\crefname{originalsection}{original section}{original sections}
\Crefname{originalsection}{Original section}{Original sections}
\makeatletter
\newlabel{stacks-morphisms-remark-identify-automorphism-groups}{{0DTW}{source}{}{sourcepointer-32}{}}
\newlabel{stacks-morphisms-remark-identify-automorphism-groups@cref}{{[originalsection][0][]0DTW}{[1][0][]source}{}{}{}}
\makeatother
\crefname{originallemma}{original lemma}{original lemmas}
\Crefname{originallemma}{Original lemma}{Original lemmas}
\makeatletter
\newlabel{stacks-morphisms-lemma-aut-iso-unramified}{{0DU9}{source}{}{sourcepointer-33}{}}
\newlabel{stacks-morphisms-lemma-aut-iso-unramified@cref}{{[originallemma][0][]0DU9}{[1][0][]source}{}{}{}}
\makeatother
\crefname{originallemma}{original lemma}{original lemmas}
\Crefname{originallemma}{Original lemma}{Original lemmas}
\makeatletter
\newlabel{stacks-morphisms-lemma-stabilizer-preserving}{{0DU7}{source}{}{sourcepointer-34}{}}
\newlabel{stacks-morphisms-lemma-stabilizer-preserving@cref}{{[originallemma][0][]0DU7}{[1][0][]source}{}{}{}}
\makeatother
\crefname{originallemma}{original lemma}{original lemmas}
\Crefname{originallemma}{Original lemma}{Original lemmas}
\makeatletter
\newlabel{groupoids-lemma-etale}{{0DT9}{source}{}{sourcepointer-35}{}}
\newlabel{groupoids-lemma-etale@cref}{{[originallemma][0][]0DT9}{[1][0][]source}{}{}{}}
\makeatother
\crefname{originallemma}{original lemma}{original lemmas}
\Crefname{originallemma}{Original lemma}{Original lemmas}
\makeatletter
\newlabel{morphisms-lemma-universal-homeomorphism}{{04DF}{source}{}{sourcepointer-36}{}}
\newlabel{morphisms-lemma-universal-homeomorphism@cref}{{[originallemma][0][]04DF}{[1][0][]source}{}{}{}}
\makeatother
\crefname{originallemma}{original lemma}{original lemmas}
\Crefname{originallemma}{Original lemma}{Original lemmas}
\makeatletter
\newlabel{stacks-morphisms-lemma-base-change-stabilizer-preserving}{{0DUB}{source}{}{sourcepointer-37}{}}
\newlabel{stacks-morphisms-lemma-base-change-stabilizer-preserving@cref}{{[originallemma][0][]0DUB}{[1][0][]source}{}{}{}}
\makeatother
\crefname{originallemma}{original lemma}{original lemmas}
\Crefname{originallemma}{Original lemma}{Original lemmas}
\makeatletter
\newlabel{spaces-properties-lemma-points-cartesian}{{03H4}{source}{}{sourcepointer-38}{}}
\newlabel{spaces-properties-lemma-points-cartesian@cref}{{[originallemma][0][]03H4}{[1][0][]source}{}{}{}}
\makeatother
\crefname{originallemma}{original lemma}{original lemmas}
\Crefname{originallemma}{Original lemma}{Original lemmas}
\makeatletter
\newlabel{more-morphisms-lemma-separated-locally-quasi-finite-over-affine}{{07S0}{source}{}{sourcepointer-39}{}}
\newlabel{more-morphisms-lemma-separated-locally-quasi-finite-over-affine@cref}{{[originallemma][0][]07S0}{[1][0][]source}{}{}{}}
\makeatother
\crefname{originallemma}{original lemma}{original lemmas}
\Crefname{originallemma}{Original lemma}{Original lemmas}
\makeatletter
\newlabel{groupoids-lemma-find-invariant-affine}{{03JE}{source}{}{sourcepointer-40}{}}
\newlabel{groupoids-lemma-find-invariant-affine@cref}{{[originallemma][0][]03JE}{[1][0][]source}{}{}{}}
\makeatother
\crefname{originallemma}{original lemma}{original lemmas}
\Crefname{originallemma}{Original lemma}{Original lemmas}
\makeatletter
\newlabel{spaces-morphisms-lemma-locally-quasi-finite-separated-representable}{{0418}{source}{}{sourcepointer-41}{}}
\newlabel{spaces-morphisms-lemma-locally-quasi-finite-separated-representable@cref}{{[originallemma][0][]0418}{[1][0][]source}{}{}{}}
\makeatother
\crefname{originalsection}{original section}{original sections}
\Crefname{originalsection}{Original section}{Original sections}
\makeatletter
\newlabel{schemes-section-glueing-schemes}{{01JA}{source}{}{sourcepointer-42}{}}
\newlabel{schemes-section-glueing-schemes@cref}{{[originalsection][0][]01JA}{[1][0][]source}{}{}{}}
\makeatother
\crefname{originallemma}{original lemma}{original lemmas}
\Crefname{originallemma}{Original lemma}{Original lemmas}
\makeatletter
\newlabel{schemes-lemma-glue}{{01JB}{source}{}{sourcepointer-43}{}}
\newlabel{schemes-lemma-glue@cref}{{[originallemma][0][]01JB}{[1][0][]source}{}{}{}}
\makeatother
\crefname{originallemma}{original lemma}{original lemmas}
\Crefname{originallemma}{Original lemma}{Original lemmas}
\makeatletter
\newlabel{morphisms-lemma-image-universally-closed-separated}{{09MQ}{source}{}{sourcepointer-44}{}}
\newlabel{morphisms-lemma-image-universally-closed-separated@cref}{{[originallemma][0][]09MQ}{[1][0][]source}{}{}{}}
\makeatother
\crefname{originallemma}{original lemma}{original lemmas}
\Crefname{originallemma}{Original lemma}{Original lemmas}
\makeatletter
\newlabel{stacks-morphisms-lemma-etale-local-quasi-DM-at-x}{{0DU4}{source}{}{sourcepointer-45}{}}
\newlabel{stacks-morphisms-lemma-etale-local-quasi-DM-at-x@cref}{{[originallemma][0][]0DU4}{[1][0][]source}{}{}{}}
\makeatother
\crefname{originallemma}{original lemma}{original lemmas}
\Crefname{originallemma}{Original lemma}{Original lemmas}
\makeatletter
\newlabel{stacks-morphisms-lemma-stabilizer-preserving-unramified}{{0DUA}{source}{}{sourcepointer-46}{}}
\newlabel{stacks-morphisms-lemma-stabilizer-preserving-unramified@cref}{{[originallemma][0][]0DUA}{[1][0][]source}{}{}{}}
\makeatother
\crefname{originallemma}{original lemma}{original lemmas}
\Crefname{originallemma}{Original lemma}{Original lemmas}
\makeatletter
\newlabel{properties-lemma-ample-finite-set-in-affine}{{01ZY}{source}{}{sourcepointer-47}{}}
\newlabel{properties-lemma-ample-finite-set-in-affine@cref}{{[originallemma][0][]01ZY}{[1][0][]source}{}{}{}}
\makeatother
\crefname{originallemma}{original lemma}{original lemmas}
\Crefname{originallemma}{Original lemma}{Original lemmas}
\makeatletter
\newlabel{stacks-morphisms-lemma-diagonal-diagonal}{{0CL0}{source}{}{sourcepointer-48}{}}
\newlabel{stacks-morphisms-lemma-diagonal-diagonal@cref}{{[originallemma][0][]0CL0}{[1][0][]source}{}{}{}}
\makeatother
\crefname{originallemma}{original lemma}{original lemmas}
\Crefname{originallemma}{Original lemma}{Original lemmas}
\makeatletter
\newlabel{stacks-morphisms-lemma-base-change-separated}{{050F}{source}{}{sourcepointer-49}{}}
\newlabel{stacks-morphisms-lemma-base-change-separated@cref}{{[originallemma][0][]050F}{[1][0][]source}{}{}{}}
\makeatother
\crefname{originallemma}{original lemma}{original lemmas}
\Crefname{originallemma}{Original lemma}{Original lemmas}
\makeatletter
\newlabel{stacks-morphisms-lemma-base-change-etale}{{0CIN}{source}{}{sourcepointer-50}{}}
\newlabel{stacks-morphisms-lemma-base-change-etale@cref}{{[originallemma][0][]0CIN}{[1][0][]source}{}{}{}}
\makeatother
\crefname{originallemma}{original lemma}{original lemmas}
\Crefname{originallemma}{Original lemma}{Original lemmas}
\makeatletter
\newlabel{stacks-morphisms-lemma-universally-injective}{{0CPN}{source}{}{sourcepointer-51}{}}
\newlabel{stacks-morphisms-lemma-universally-injective@cref}{{[originallemma][0][]0CPN}{[1][0][]source}{}{}{}}
\makeatother
\crefname{originallemma}{original lemma}{original lemmas}
\Crefname{originallemma}{Original lemma}{Original lemmas}
\makeatletter
\newlabel{stacks-morphisms-lemma-universally-injective-point}{{0DTM}{source}{}{sourcepointer-52}{}}
\newlabel{stacks-morphisms-lemma-universally-injective-point@cref}{{[originallemma][0][]0DTM}{[1][0][]source}{}{}{}}
\makeatother
\crefname{originallemma}{original lemma}{original lemmas}
\Crefname{originallemma}{Original lemma}{Original lemmas}
\makeatletter
\newlabel{groupoids-lemma-etale-equivalence-relation}{{0DT7}{source}{}{sourcepointer-53}{}}
\newlabel{groupoids-lemma-etale-equivalence-relation@cref}{{[originallemma][0][]0DT7}{[1][0][]source}{}{}{}}
\makeatother
\crefname{originaltheorem}{original theorem}{original theorems}
\Crefname{originaltheorem}{Original theorem}{Original theorems}
\makeatletter
\newlabel{spaces-theorem-presentation}{{02WW}{source}{}{sourcepointer-54}{}}
\newlabel{spaces-theorem-presentation@cref}{{[originaltheorem][0][]02WW}{[1][0][]source}{}{}{}}
\makeatother
\crefname{originallemma}{original lemma}{original lemmas}
\Crefname{originallemma}{Original lemma}{Original lemmas}
\makeatletter
\newlabel{stacks-morphisms-lemma-etale-permanence}{{0CIR}{source}{}{sourcepointer-55}{}}
\newlabel{stacks-morphisms-lemma-etale-permanence@cref}{{[originallemma][0][]0CIR}{[1][0][]source}{}{}{}}
\makeatother
\crefname{originallemma}{original lemma}{original lemmas}
\Crefname{originallemma}{Original lemma}{Original lemmas}
\makeatletter
\newlabel{stacks-morphisms-lemma-etale-iso}{{0DUD}{source}{}{sourcepointer-56}{}}
\newlabel{stacks-morphisms-lemma-etale-iso@cref}{{[originallemma][0][]0DUD}{[1][0][]source}{}{}{}}
\makeatother
\crefname{originallemma}{original lemma}{original lemmas}
\Crefname{originallemma}{Original lemma}{Original lemmas}
\makeatletter
\newlabel{stacks-morphisms-lemma-check-separated-covering}{{06TZ}{source}{}{sourcepointer-57}{}}
\newlabel{stacks-morphisms-lemma-check-separated-covering@cref}{{[originallemma][0][]06TZ}{[1][0][]source}{}{}{}}
\makeatother
\crefname{originallemma}{original lemma}{original lemmas}
\Crefname{originallemma}{Original lemma}{Original lemmas}
\makeatletter
\newlabel{stacks-morphisms-lemma-check-quasi-compact-covering}{{0DTL}{source}{}{sourcepointer-58}{}}
\newlabel{stacks-morphisms-lemma-check-quasi-compact-covering@cref}{{[originallemma][0][]0DTL}{[1][0][]source}{}{}{}}
\makeatother
\crefname{originallemma}{original lemma}{original lemmas}
\Crefname{originallemma}{Original lemma}{Original lemmas}
\makeatletter
\newlabel{stacks-morphisms-lemma-check-universal-homeomorphism-covering}{{0DTQ}{source}{}{sourcepointer-59}{}}
\newlabel{stacks-morphisms-lemma-check-universal-homeomorphism-covering@cref}{{[originallemma][0][]0DTQ}{[1][0][]source}{}{}{}}
\makeatother
\title{Keel–Mori uniform categorical moduli space for algebraic stacks with finite inertia}
\author{The Stacks Project Authors}
\date{stacks-project:stacks-more-morphisms.tex:local-snapshot-20261011; GNU FDL 1.2 or later}
\begin{document}
\maketitle
\noindent\textit{Editorial note.} The two prior substantive reviews are reused for the global proof and its compressed functoriality checks. This preparation makes no new claim that every upstream invariant-theory theorem was re-proved.\par\medskip
\noindent\textit{Editorial note.} The author omits routine verification of base-change functoriality and commutativity at3782–3785; the local definitions and etale sheaf gluing supply that verification.\par\medskip
\noindent\textit{Allowed standard prerequisites (editorial).}\begin{itemize}
\item Explicit prerequisite boundary. Algebraic stacks and inertia, etale descent and quotient algebraic spaces for etale equivalence relations, finite locally free groupoid invariant quotients, and descent of separated/quasi-compact/universal-homeomorphism properties. Local finite-inertia presentations and moduli-space lemmas are retained.
\end{itemize}\medskip
\ifdefined\hypertarget\hypertarget{sourcepointer-1}{}\fi
\noindent\textit{Source reference.} Original section 048B: focused-seeds/AGP/sources/stacks-snapshot-20261011/tags.txt, source lines 5538--5538. Stable Stacks tag for an expressly allowed foundational or notational input; the counted construction remains the two concrete core obligations. The retained local core is not replaced by this pointer. See standard\_prerequisites.\par\medskip
\ifdefined\hypertarget\hypertarget{sourcepointer-2}{}\fi
\noindent\textit{Source reference.} Original theorem 06FI: focused-seeds/AGP/sources/stacks-snapshot-20261011/tags.txt, source lines 8252--8252. Stable Stacks tag for an expressly allowed foundational or notational input; the counted construction remains the two concrete core obligations. The retained local core is not replaced by this pointer. See standard\_prerequisites.\par\medskip
\ifdefined\hypertarget\hypertarget{sourcepointer-3}{}\fi
\noindent\textit{Source reference.} Original lemma 04M9: focused-seeds/AGP/sources/stacks-snapshot-20261011/tags.txt, source lines 6035--6035. Stable Stacks tag for an expressly allowed foundational or notational input; the counted construction remains the two concrete core obligations. The retained local core is not replaced by this pointer. See standard\_prerequisites.\par\medskip
\ifdefined\hypertarget\hypertarget{sourcepointer-4}{}\fi
\noindent\textit{Source reference.} Original lemma 044U: focused-seeds/AGP/sources/stacks-snapshot-20261011/tags.txt, source lines 5413--5413. Stable Stacks tag for an expressly allowed foundational or notational input; the counted construction remains the two concrete core obligations. The retained local core is not replaced by this pointer. See standard\_prerequisites.\par\medskip
\ifdefined\hypertarget\hypertarget{sourcepointer-5}{}\fi
\noindent\textit{Source reference.} Original lemma 044X: focused-seeds/AGP/sources/stacks-snapshot-20261011/tags.txt, source lines 5416--5416. Stable Stacks tag for an expressly allowed foundational or notational input; the counted construction remains the two concrete core obligations. The retained local core is not replaced by this pointer. See standard\_prerequisites.\par\medskip
\ifdefined\hypertarget\hypertarget{sourcepointer-6}{}\fi
\noindent\textit{Source reference.} Original definition 048J: focused-seeds/AGP/sources/stacks-snapshot-20261011/tags.txt, source lines 5546--5546. Stable Stacks tag for an expressly allowed foundational or notational input; the counted construction remains the two concrete core obligations. The retained local core is not replaced by this pointer. See standard\_prerequisites.\par\medskip
\ifdefined\hypertarget\hypertarget{sourcepointer-7}{}\fi
\noindent\textit{Source reference.} Original lemma 06FH: focused-seeds/AGP/sources/stacks-snapshot-20261011/tags.txt, source lines 8251--8251. Stable Stacks tag for an expressly allowed foundational or notational input; the counted construction remains the two concrete core obligations. The retained local core is not replaced by this pointer. See standard\_prerequisites.\par\medskip
\ifdefined\hypertarget\hypertarget{sourcepointer-8}{}\fi
\noindent\textit{Source reference.} Original section 04WY: focused-seeds/AGP/sources/stacks-snapshot-20261011/tags.txt, source lines 6376--6376. Stable Stacks tag for an expressly allowed foundational or notational input; the counted construction remains the two concrete core obligations. The retained local core is not replaced by this pointer. See standard\_prerequisites.\par\medskip
\ifdefined\hypertarget\hypertarget{sourcepointer-9}{}\fi
\noindent\textit{Source reference.} Original section 0261: focused-seeds/AGP/sources/stacks-snapshot-20261011/tags.txt, source lines 2834--2834. Stable Stacks tag for an expressly allowed foundational or notational input; the counted construction remains the two concrete core obligations. The retained local core is not replaced by this pointer. See standard\_prerequisites.\par\medskip
\ifdefined\hypertarget\hypertarget{sourcepointer-10}{}\fi
\noindent\textit{Source reference.} Original section 04UX: focused-seeds/AGP/sources/stacks-snapshot-20261011/tags.txt, source lines 6305--6305. Stable Stacks tag for an expressly allowed foundational or notational input; the counted construction remains the two concrete core obligations. The retained local core is not replaced by this pointer. See standard\_prerequisites.\par\medskip
\ifdefined\hypertarget\hypertarget{sourcepointer-11}{}\fi
\noindent\textit{Source reference.} Original lemma 04T4: focused-seeds/AGP/sources/stacks-snapshot-20261011/tags.txt, source lines 6242--6242. Stable Stacks tag for an expressly allowed foundational or notational input; the counted construction remains the two concrete core obligations. The retained local core is not replaced by this pointer. See standard\_prerequisites.\par\medskip
\ifdefined\hypertarget\hypertarget{sourcepointer-12}{}\fi
\noindent\textit{Source reference.} Original lemma 02KB: focused-seeds/AGP/sources/stacks-snapshot-20261011/tags.txt, source lines 3366--3366. Stable Stacks tag for an expressly allowed foundational or notational input; the counted construction remains the two concrete core obligations. The retained local core is not replaced by this pointer. See standard\_prerequisites.\par\medskip
\ifdefined\hypertarget\hypertarget{sourcepointer-13}{}\fi
\noindent\textit{Source reference.} Original lemma 04XD: focused-seeds/AGP/sources/stacks-snapshot-20261011/tags.txt, source lines 6391--6391. Stable Stacks tag for an expressly allowed foundational or notational input; the counted construction remains the two concrete core obligations. The retained local core is not replaced by this pointer. See standard\_prerequisites.\par\medskip
\ifdefined\hypertarget\hypertarget{sourcepointer-14}{}\fi
\noindent\textit{Source reference.} Original lemma 04YC: focused-seeds/AGP/sources/stacks-snapshot-20261011/tags.txt, source lines 6425--6425. Stable Stacks tag for an expressly allowed foundational or notational input; the counted construction remains the two concrete core obligations. The retained local core is not replaced by this pointer. See standard\_prerequisites.\par\medskip
\ifdefined\hypertarget\hypertarget{sourcepointer-15}{}\fi
\noindent\textit{Source reference.} Original lemma 0DTY: focused-seeds/AGP/sources/stacks-snapshot-20261011/tags.txt, source lines 17308--17308. Stable Stacks tag for an expressly allowed foundational or notational input; the counted construction remains the two concrete core obligations. The retained local core is not replaced by this pointer. See standard\_prerequisites.\par\medskip
\ifdefined\hypertarget\hypertarget{sourcepointer-16}{}\fi
\noindent\textit{Source reference.} Original lemma 07VS: focused-seeds/AGP/sources/stacks-snapshot-20261011/tags.txt, source lines 10017--10017. Stable Stacks tag for an expressly allowed foundational or notational input; the counted construction remains the two concrete core obligations. The retained local core is not replaced by this pointer. See standard\_prerequisites.\par\medskip
\ifdefined\hypertarget\hypertarget{sourcepointer-17}{}\fi
\noindent\textit{Source reference.} Original proposition 03BM: focused-seeds/AGP/sources/stacks-snapshot-20261011/tags.txt, source lines 4353--4353. Stable Stacks tag for an expressly allowed foundational or notational input; the counted construction remains the two concrete core obligations. The retained local core is not replaced by this pointer. See standard\_prerequisites.\par\medskip
\ifdefined\hypertarget\hypertarget{sourcepointer-18}{}\fi
\noindent\textit{Source reference.} Original lemma 0CHR: focused-seeds/AGP/sources/stacks-snapshot-20261011/tags.txt, source lines 15691--15691. Stable Stacks tag for an expressly allowed foundational or notational input; the counted construction remains the two concrete core obligations. The retained local core is not replaced by this pointer. See standard\_prerequisites.\par\medskip
\ifdefined\hypertarget\hypertarget{sourcepointer-19}{}\fi
\noindent\textit{Source reference.} Original section 03BE: focused-seeds/AGP/sources/stacks-snapshot-20261011/tags.txt, source lines 4345--4345. Stable Stacks tag for an expressly allowed foundational or notational input; the counted construction remains the two concrete core obligations. The retained local core is not replaced by this pointer. See standard\_prerequisites.\par\medskip
\ifdefined\hypertarget\hypertarget{sourcepointer-20}{}\fi
\noindent\textit{Source reference.} Original lemma 01I1: focused-seeds/AGP/sources/stacks-snapshot-20261011/tags.txt, source lines 1962--1962. Stable Stacks tag for an expressly allowed foundational or notational input; the counted construction remains the two concrete core obligations. The retained local core is not replaced by this pointer. See standard\_prerequisites.\par\medskip
\ifdefined\hypertarget\hypertarget{sourcepointer-21}{}\fi
\noindent\textit{Source reference.} Original lemma 050M: focused-seeds/AGP/sources/stacks-snapshot-20261011/tags.txt, source lines 6505--6505. Stable Stacks tag for an expressly allowed foundational or notational input; the counted construction remains the two concrete core obligations. The retained local core is not replaced by this pointer. See standard\_prerequisites.\par\medskip
\ifdefined\hypertarget\hypertarget{sourcepointer-22}{}\fi
\noindent\textit{Source reference.} Original lemma 050X: focused-seeds/AGP/sources/stacks-snapshot-20261011/tags.txt, source lines 6515--6515. Stable Stacks tag for an expressly allowed foundational or notational input; the counted construction remains the two concrete core obligations. The retained local core is not replaced by this pointer. See standard\_prerequisites.\par\medskip
\ifdefined\hypertarget\hypertarget{sourcepointer-23}{}\fi
\noindent\textit{Source reference.} Original lemma 03BJ: focused-seeds/AGP/sources/stacks-snapshot-20261011/tags.txt, source lines 4350--4350. Stable Stacks tag for an expressly allowed foundational or notational input; the counted construction remains the two concrete core obligations. The retained local core is not replaced by this pointer. See standard\_prerequisites.\par\medskip
\ifdefined\hypertarget\hypertarget{sourcepointer-24}{}\fi
\noindent\textit{Source reference.} Original lemma 01WM: focused-seeds/AGP/sources/stacks-snapshot-20261011/tags.txt, source lines 2487--2487. Stable Stacks tag for an expressly allowed foundational or notational input; the counted construction remains the two concrete core obligations. The retained local core is not replaced by this pointer. See standard\_prerequisites.\par\medskip
\ifdefined\hypertarget\hypertarget{sourcepointer-25}{}\fi
\noindent\textit{Source reference.} Original lemma 0CQK: focused-seeds/AGP/sources/stacks-snapshot-20261011/tags.txt, source lines 15965--15965. Stable Stacks tag for an expressly allowed foundational or notational input; the counted construction remains the two concrete core obligations. The retained local core is not replaced by this pointer. See standard\_prerequisites.\par\medskip
\ifdefined\hypertarget\hypertarget{sourcepointer-26}{}\fi
\noindent\textit{Source reference.} Original lemma 0DTZ: focused-seeds/AGP/sources/stacks-snapshot-20261011/tags.txt, source lines 17309--17309. Stable Stacks tag for an expressly allowed foundational or notational input; the counted construction remains the two concrete core obligations. The retained local core is not replaced by this pointer. See standard\_prerequisites.\par\medskip
\ifdefined\hypertarget\hypertarget{sourcepointer-27}{}\fi
\noindent\textit{Source reference.} Original lemma 03BL: focused-seeds/AGP/sources/stacks-snapshot-20261011/tags.txt, source lines 4352--4352. Stable Stacks tag for an expressly allowed foundational or notational input; the counted construction remains the two concrete core obligations. The retained local core is not replaced by this pointer. See standard\_prerequisites.\par\medskip
\ifdefined\hypertarget\hypertarget{sourcepointer-28}{}\fi
\noindent\textit{Source reference.} Original lemma 0DTG: focused-seeds/AGP/sources/stacks-snapshot-20261011/tags.txt, source lines 17291--17291. Stable Stacks tag for an expressly allowed foundational or notational input; the counted construction remains the two concrete core obligations. The retained local core is not replaced by this pointer. See standard\_prerequisites.\par\medskip
\ifdefined\hypertarget\hypertarget{sourcepointer-29}{}\fi
\noindent\textit{Source reference.} Original lemma 03BK: focused-seeds/AGP/sources/stacks-snapshot-20261011/tags.txt, source lines 4351--4351. Stable Stacks tag for an expressly allowed foundational or notational input; the counted construction remains the two concrete core obligations. The retained local core is not replaced by this pointer. See standard\_prerequisites.\par\medskip
\ifdefined\hypertarget\hypertarget{sourcepointer-30}{}\fi
\noindent\textit{Source reference.} Original lemma 09EF: focused-seeds/AGP/sources/stacks-snapshot-20261011/tags.txt, source lines 11898--11898. Stable Stacks tag for an expressly allowed foundational or notational input; the counted construction remains the two concrete core obligations. The retained local core is not replaced by this pointer. See standard\_prerequisites.\par\medskip
\ifdefined\hypertarget\hypertarget{sourcepointer-31}{}\fi
\noindent\textit{Source reference.} Original lemma 0DTN: focused-seeds/AGP/sources/stacks-snapshot-20261011/tags.txt, source lines 17298--17298. Stable Stacks tag for an expressly allowed foundational or notational input; the counted construction remains the two concrete core obligations. The retained local core is not replaced by this pointer. See standard\_prerequisites.\par\medskip
\ifdefined\hypertarget\hypertarget{sourcepointer-32}{}\fi
\noindent\textit{Source reference.} Original section 0DTW: focused-seeds/AGP/sources/stacks-snapshot-20261011/tags.txt, source lines 17306--17306. Stable Stacks tag for an expressly allowed foundational or notational input; the counted construction remains the two concrete core obligations. The retained local core is not replaced by this pointer. See standard\_prerequisites.\par\medskip
\ifdefined\hypertarget\hypertarget{sourcepointer-33}{}\fi
\noindent\textit{Source reference.} Original lemma 0DU9: focused-seeds/AGP/sources/stacks-snapshot-20261011/tags.txt, source lines 17319--17319. Stable Stacks tag for an expressly allowed foundational or notational input; the counted construction remains the two concrete core obligations. The retained local core is not replaced by this pointer. See standard\_prerequisites.\par\medskip
\ifdefined\hypertarget\hypertarget{sourcepointer-34}{}\fi
\noindent\textit{Source reference.} Original lemma 0DU7: focused-seeds/AGP/sources/stacks-snapshot-20261011/tags.txt, source lines 17317--17317. Stable Stacks tag for an expressly allowed foundational or notational input; the counted construction remains the two concrete core obligations. The retained local core is not replaced by this pointer. See standard\_prerequisites.\par\medskip
\ifdefined\hypertarget\hypertarget{sourcepointer-35}{}\fi
\noindent\textit{Source reference.} Original lemma 0DT9: focused-seeds/AGP/sources/stacks-snapshot-20261011/tags.txt, source lines 17284--17284. Stable Stacks tag for an expressly allowed foundational or notational input; the counted construction remains the two concrete core obligations. The retained local core is not replaced by this pointer. See standard\_prerequisites.\par\medskip
\ifdefined\hypertarget\hypertarget{sourcepointer-36}{}\fi
\noindent\textit{Source reference.} Original lemma 04DF: focused-seeds/AGP/sources/stacks-snapshot-20261011/tags.txt, source lines 5725--5725. Stable Stacks tag for an expressly allowed foundational or notational input; the counted construction remains the two concrete core obligations. The retained local core is not replaced by this pointer. See standard\_prerequisites.\par\medskip
\ifdefined\hypertarget\hypertarget{sourcepointer-37}{}\fi
\noindent\textit{Source reference.} Original lemma 0DUB: focused-seeds/AGP/sources/stacks-snapshot-20261011/tags.txt, source lines 17321--17321. Stable Stacks tag for an expressly allowed foundational or notational input; the counted construction remains the two concrete core obligations. The retained local core is not replaced by this pointer. See standard\_prerequisites.\par\medskip
\ifdefined\hypertarget\hypertarget{sourcepointer-38}{}\fi
\noindent\textit{Source reference.} Original lemma 03H4: focused-seeds/AGP/sources/stacks-snapshot-20261011/tags.txt, source lines 4553--4553. Stable Stacks tag for an expressly allowed foundational or notational input; the counted construction remains the two concrete core obligations. The retained local core is not replaced by this pointer. See standard\_prerequisites.\par\medskip
\ifdefined\hypertarget\hypertarget{sourcepointer-39}{}\fi
\noindent\textit{Source reference.} Original lemma 07S0: focused-seeds/AGP/sources/stacks-snapshot-20261011/tags.txt, source lines 9885--9885. Stable Stacks tag for an expressly allowed foundational or notational input; the counted construction remains the two concrete core obligations. The retained local core is not replaced by this pointer. See standard\_prerequisites.\par\medskip
\ifdefined\hypertarget\hypertarget{sourcepointer-40}{}\fi
\noindent\textit{Source reference.} Original lemma 03JE: focused-seeds/AGP/sources/stacks-snapshot-20261011/tags.txt, source lines 4636--4636. Stable Stacks tag for an expressly allowed foundational or notational input; the counted construction remains the two concrete core obligations. The retained local core is not replaced by this pointer. See standard\_prerequisites.\par\medskip
\ifdefined\hypertarget\hypertarget{sourcepointer-41}{}\fi
\noindent\textit{Source reference.} Original lemma 0418: focused-seeds/AGP/sources/stacks-snapshot-20261011/tags.txt, source lines 5283--5283. Stable Stacks tag for an expressly allowed foundational or notational input; the counted construction remains the two concrete core obligations. The retained local core is not replaced by this pointer. See standard\_prerequisites.\par\medskip
\ifdefined\hypertarget\hypertarget{sourcepointer-42}{}\fi
\noindent\textit{Source reference.} Original section 01JA: focused-seeds/AGP/sources/stacks-snapshot-20261011/tags.txt, source lines 2007--2007. Stable Stacks tag for an expressly allowed foundational or notational input; the counted construction remains the two concrete core obligations. The retained local core is not replaced by this pointer. See standard\_prerequisites.\par\medskip
\ifdefined\hypertarget\hypertarget{sourcepointer-43}{}\fi
\noindent\textit{Source reference.} Original lemma 01JB: focused-seeds/AGP/sources/stacks-snapshot-20261011/tags.txt, source lines 2008--2008. Stable Stacks tag for an expressly allowed foundational or notational input; the counted construction remains the two concrete core obligations. The retained local core is not replaced by this pointer. See standard\_prerequisites.\par\medskip
\ifdefined\hypertarget\hypertarget{sourcepointer-44}{}\fi
\noindent\textit{Source reference.} Original lemma 09MQ: focused-seeds/AGP/sources/stacks-snapshot-20261011/tags.txt, source lines 12188--12188. Stable Stacks tag for an expressly allowed foundational or notational input; the counted construction remains the two concrete core obligations. The retained local core is not replaced by this pointer. See standard\_prerequisites.\par\medskip
\ifdefined\hypertarget\hypertarget{sourcepointer-45}{}\fi
\noindent\textit{Source reference.} Original lemma 0DU4: focused-seeds/AGP/sources/stacks-snapshot-20261011/tags.txt, source lines 17314--17314. Stable Stacks tag for an expressly allowed foundational or notational input; the counted construction remains the two concrete core obligations. The retained local core is not replaced by this pointer. See standard\_prerequisites.\par\medskip
\ifdefined\hypertarget\hypertarget{sourcepointer-46}{}\fi
\noindent\textit{Source reference.} Original lemma 0DUA: focused-seeds/AGP/sources/stacks-snapshot-20261011/tags.txt, source lines 17320--17320. Stable Stacks tag for an expressly allowed foundational or notational input; the counted construction remains the two concrete core obligations. The retained local core is not replaced by this pointer. See standard\_prerequisites.\par\medskip
\ifdefined\hypertarget\hypertarget{sourcepointer-47}{}\fi
\noindent\textit{Source reference.} Original lemma 01ZY: focused-seeds/AGP/sources/stacks-snapshot-20261011/tags.txt, source lines 2613--2613. Stable Stacks tag for an expressly allowed foundational or notational input; the counted construction remains the two concrete core obligations. The retained local core is not replaced by this pointer. See standard\_prerequisites.\par\medskip
\ifdefined\hypertarget\hypertarget{sourcepointer-48}{}\fi
\noindent\textit{Source reference.} Original lemma 0CL0: focused-seeds/AGP/sources/stacks-snapshot-20261011/tags.txt, source lines 15805--15805. Stable Stacks tag for an expressly allowed foundational or notational input; the counted construction remains the two concrete core obligations. The retained local core is not replaced by this pointer. See standard\_prerequisites.\par\medskip
\ifdefined\hypertarget\hypertarget{sourcepointer-49}{}\fi
\noindent\textit{Source reference.} Original lemma 050F: focused-seeds/AGP/sources/stacks-snapshot-20261011/tags.txt, source lines 6498--6498. Stable Stacks tag for an expressly allowed foundational or notational input; the counted construction remains the two concrete core obligations. The retained local core is not replaced by this pointer. See standard\_prerequisites.\par\medskip
\ifdefined\hypertarget\hypertarget{sourcepointer-50}{}\fi
\noindent\textit{Source reference.} Original lemma 0CIN: focused-seeds/AGP/sources/stacks-snapshot-20261011/tags.txt, source lines 15723--15723. Stable Stacks tag for an expressly allowed foundational or notational input; the counted construction remains the two concrete core obligations. The retained local core is not replaced by this pointer. See standard\_prerequisites.\par\medskip
\ifdefined\hypertarget\hypertarget{sourcepointer-51}{}\fi
\noindent\textit{Source reference.} Original lemma 0CPN: focused-seeds/AGP/sources/stacks-snapshot-20261011/tags.txt, source lines 15933--15933. Stable Stacks tag for an expressly allowed foundational or notational input; the counted construction remains the two concrete core obligations. The retained local core is not replaced by this pointer. See standard\_prerequisites.\par\medskip
\ifdefined\hypertarget\hypertarget{sourcepointer-52}{}\fi
\noindent\textit{Source reference.} Original lemma 0DTM: focused-seeds/AGP/sources/stacks-snapshot-20261011/tags.txt, source lines 17297--17297. Stable Stacks tag for an expressly allowed foundational or notational input; the counted construction remains the two concrete core obligations. The retained local core is not replaced by this pointer. See standard\_prerequisites.\par\medskip
\ifdefined\hypertarget\hypertarget{sourcepointer-53}{}\fi
\noindent\textit{Source reference.} Original lemma 0DT7: focused-seeds/AGP/sources/stacks-snapshot-20261011/tags.txt, source lines 17282--17282. Stable Stacks tag for an expressly allowed foundational or notational input; the counted construction remains the two concrete core obligations. The retained local core is not replaced by this pointer. See standard\_prerequisites.\par\medskip
\ifdefined\hypertarget\hypertarget{sourcepointer-54}{}\fi
\noindent\textit{Source reference.} Original theorem 02WW: focused-seeds/AGP/sources/stacks-snapshot-20261011/tags.txt, source lines 3819--3819. Stable Stacks tag for an expressly allowed foundational or notational input; the counted construction remains the two concrete core obligations. The retained local core is not replaced by this pointer. See standard\_prerequisites.\par\medskip
\ifdefined\hypertarget\hypertarget{sourcepointer-55}{}\fi
\noindent\textit{Source reference.} Original lemma 0CIR: focused-seeds/AGP/sources/stacks-snapshot-20261011/tags.txt, source lines 15726--15726. Stable Stacks tag for an expressly allowed foundational or notational input; the counted construction remains the two concrete core obligations. The retained local core is not replaced by this pointer. See standard\_prerequisites.\par\medskip
\ifdefined\hypertarget\hypertarget{sourcepointer-56}{}\fi
\noindent\textit{Source reference.} Original lemma 0DUD: focused-seeds/AGP/sources/stacks-snapshot-20261011/tags.txt, source lines 17323--17323. Stable Stacks tag for an expressly allowed foundational or notational input; the counted construction remains the two concrete core obligations. The retained local core is not replaced by this pointer. See standard\_prerequisites.\par\medskip
\ifdefined\hypertarget\hypertarget{sourcepointer-57}{}\fi
\noindent\textit{Source reference.} Original lemma 06TZ: focused-seeds/AGP/sources/stacks-snapshot-20261011/tags.txt, source lines 8728--8728. Stable Stacks tag for an expressly allowed foundational or notational input; the counted construction remains the two concrete core obligations. The retained local core is not replaced by this pointer. See standard\_prerequisites.\par\medskip
\ifdefined\hypertarget\hypertarget{sourcepointer-58}{}\fi
\noindent\textit{Source reference.} Original lemma 0DTL: focused-seeds/AGP/sources/stacks-snapshot-20261011/tags.txt, source lines 17296--17296. Stable Stacks tag for an expressly allowed foundational or notational input; the counted construction remains the two concrete core obligations. The retained local core is not replaced by this pointer. See standard\_prerequisites.\par\medskip
\ifdefined\hypertarget\hypertarget{sourcepointer-59}{}\fi
\noindent\textit{Source reference.} Original lemma 0DTQ: focused-seeds/AGP/sources/stacks-snapshot-20261011/tags.txt, source lines 17300--17300. Stable Stacks tag for an expressly allowed foundational or notational input; the counted construction remains the two concrete core obligations. The retained local core is not replaced by this pointer. See standard\_prerequisites.\par\medskip
\begin{theorem}[Keel-Mori]
\label{theorem-keel-mori}
Let $\mathcal{X}$ be an algebraic stack. Assume
$\mathcal{I}_\mathcal{X} \to \mathcal{X}$ is finite.
Then there exists a uniform categorical moduli space
$$
f : \mathcal{X} \longrightarrow M
$$
and $f$ is separated, quasi-compact, and a universal homeomorphism.
\end{theorem}

\section{Moduli spaces}
\label{section-moduli-spaces}

\noindent
This section discusses morphisms $f : \mathcal{X} \to Y$ from
algebraic stacks to algebraic spaces. Under suitable hypotheses
$Y$ is called a {\it moduli space} for $\mathcal{X}$. If
$\mathcal{X} = [U/R]$ is a presentation, then we obtain an
$R$-invariant morphism $U \to Y$ and (under suitable hypotheses)
$Y$ is a {\it quotient} of the groupoid $(U, R, s, t, c)$.
A discussion of the different types of quotients can be found
starting with
Quotients of Groupoids, Section \ref{groupoids-quotients-section-introduction}.

\begin{definition}
\label{definition-categorical-quotient}
Let $\mathcal{X}$ be an algebraic stack. Let
$f : \mathcal{X} \to Y$ be a morphism to an algebraic space $Y$.
\begin{enumerate}
\item We say $f$ is a {\it categorical moduli space} if any morphism
$\mathcal{X} \to W$ to an algebraic space $W$ factors uniquely through $f$.
\item We say $f$ is a {\it uniform categorical moduli space}
if for any flat morphism $Y' \to Y$ of algebraic spaces the base change
$f' : Y' \times_Y \mathcal{X} \to Y'$ is a categorical moduli space.
\end{enumerate}
Let $\mathcal{C}$ be a full subcategory of the category of algebraic
spaces.
\begin{enumerate}
\item[(3)] We say $f$ is a {\it categorical moduli space in $\mathcal{C}$}
if $Y \in \Ob(\mathcal{C})$ and any morphism $\mathcal{X} \to W$ with
$W \in \Ob(\mathcal{C})$ factors uniquely through $f$.
\item[(4)] We say is a {\it uniform categorical moduli space in $\mathcal{C}$}
if $Y \in \Ob(\mathcal{C})$ and for every flat morphism $Y' \to Y$ in
$\mathcal{C}$ the base change $f' : Y' \times_Y \mathcal{X} \to Y'$ is a
categorical moduli space in $\mathcal{C}$.
\end{enumerate}
\end{definition}

\noindent
By the Yoneda lemma a categorical moduli space, if it exists, is unique.
Let us match this with the language introduced for quotients.

\begin{lemma}
\label{lemma-quotient-compare}
Let $(U, R, s, t, c)$ be a groupoid in algebraic spaces with
$s, t : R \to U$ flat and locally of finite presentation.
Consider the algebraic stack $\mathcal{X} = [U/R]$.
Given an algebraic space $Y$ there is a $1$-to-$1$ correspondence between
morphisms $f : \mathcal{X} \to Y$ and $R$-invariant morphisms
$\phi : U \to Y$.
\end{lemma}

\begin{proof}
Criteria for Representability, Theorem
\ref{criteria-theorem-flat-groupoid-gives-algebraic-stack}
tells us $\mathcal{X}$ is an algebraic stack.
Given a morphism $f : \mathcal{X} \to Y$ we let $\phi : U \to Y$ be
the composition $U \to \mathcal{X} \to Y$. Since $R = U \times_\mathcal{X} U$
(Groupoids in Spaces, Lemma
\ref{spaces-groupoids-lemma-quotient-stack-2-cartesian})
it is immediate that $\phi$ is $R$-invariant.
Conversely, if $\phi : U \to Y$ is an $R$-invariant morphism towards
an algebraic space, we obtain a morphism
$f : \mathcal{X} \to Y$ by
Groupoids in Spaces, Lemma
\ref{spaces-groupoids-lemma-quotient-stack-2-coequalizer}.
You can also construct $f$ from $\phi$ using the explicit description of
the quotient stack in
Groupoids in Spaces, Lemma
\ref{spaces-groupoids-lemma-quotient-stack-objects}.
\end{proof}

\begin{lemma}
\label{lemma-categorical-quotient-compare}
With assumption and notation as in Lemma \ref{lemma-quotient-compare}.
Then $f$ is a (uniform) categorical moduli space
if and only if $\phi$ is a (uniform) categorical quotient.
Similarly for moduli spaces in a full subcategory.
\end{lemma}

\begin{proof}
It is immediate from the $1$-to-$1$ correspondence established in
Lemma \ref{lemma-quotient-compare} that $f$ is a categorical moduli space
if and only if $\phi$ is a categorical quotient
(Quotients of Groupoids, Definition
\ref{groupoids-quotients-definition-categorical}).
If $Y' \to Y$ is a morphism, then
$U' = Y' \times_Y U \to Y' \times_Y \mathcal{X} = \mathcal{X}'$
is a surjective, flat, locally finitely presented morphism
as a base change of $U \to \mathcal{X}$
(Criteria for Representability, Lemma
\ref{criteria-lemma-flat-quotient-flat-presentation}).
And $R' = Y' \times_Y R$ is equal to $U' \times_{\mathcal{X}'} U'$
by transitivity of fibre products.
Hence $\mathcal{X}' = [U'/R']$, see
Algebraic Stacks, Remark \ref{algebraic-remark-flat-fp-presentation}.
Thus the base change of our situation to $Y'$ is another situation
as in the statement of the lemma. From this it immediately
follows that $f$ is a uniform categorical moduli space
if and only if $\phi$ is a uniform categorical quotient.
\end{proof}

\begin{lemma}
\label{lemma-check-uniform-categorical-quotient-on-affines}
Let $f : \mathcal{X} \to Y$ be a morphism from an algebraic stack
to an algebraic space. If for every affine scheme $Y'$ and flat
morphism $Y' \to Y$ the base change
$f' : Y' \times_Y \mathcal{X} \to Y'$ is a categorical moduli space,
then $f$ is a uniform categorical moduli space.
\end{lemma}

\begin{proof}
Choose an \'etale covering $\{Y_i \to Y\}$ where $Y_i$ is an affine scheme.
For each $i$ and $j$ choose a affine open covering
$Y_i \times_Y Y_j = \bigcup Y_{ijk}$.
Set $\mathcal{X}_i = Y_i \times_Y \mathcal{X}$ and
$\mathcal{X}_{ijk} = Y_{ijk} \times_Y \mathcal{X}$.
Let $g : \mathcal{X} \to W$ be a morphism towards
an algebraic space. Then we consider the diagram
$$
\xymatrix{
\mathcal{X}_i \ar[r] \ar[d] & \mathcal{X} \ar[d] \ar[r]_g & W \\
Y_i \ar[r] \ar@{..>}[rru] & Y
}
$$
The assumption that $\mathcal{X}_i \to Y_i$ is a categorical moduli
space, produces a unique dotted arrow $h_i : Y_i \to W$.
The assumption that $\mathcal{X}_{ijk} \to Y_{ijk}$ is a categorical
moduli space, implies the restriction of $h_i$ and $h_j$ to
$Y_{ijk}$ are equal. Hence $h_i$ and $h_j$ agree on $Y_i \times_Y Y_j$.
Since $Y = \coprod Y_i / \coprod Y_i \times_Y Y_j$
(by Spaces, Section \ref{spaces-section-presentations}) we conclude
that there is a unique morphism $Y \to W$ through which $g$ factors.
Thus $f$ is a categorical moduli space. The same argument applies
after a flat base change, hence $f$ is a uniform categorical moduli space.
\end{proof}







\section{The Keel-Mori theorem}
\label{section-Keel-Mori}

\noindent
In this section we start discussing the theorem of Keel and Mori
in the setting of algebraic stacks. For a discussion of the
literature, please see Guide to Literature, Subsection
\ref{guide-subsection-coarse-moduli-spaces}.

\begin{definition}
\label{definition-well-nigh-affine}
Let $\mathcal{X}$ be an algebraic stack. We say $\mathcal{X}$
is {\it well-nigh affine} if there exists an affine scheme $U$
and a surjective, flat, finite, and finitely presented morphism
$U \to \mathcal{X}$.
\end{definition}

\noindent
We give this property a somewhat ridiculous name because we do not intend to
use it too much.

\begin{lemma}
\label{lemma-well-nigh-affine}
Let $\mathcal{X}$ be an algebraic stack. The following are equivalent
\begin{enumerate}
\item $\mathcal{X}$ is well-nigh affine, and
\item there exists a groupoid scheme $(U, R, s, t, c)$ with $U$ and
$R$ affine and $s, t : R \to U$ finite locally free such that
$\mathcal{X} = [U/R]$.
\end{enumerate}
If true then $\mathcal{X}$ is quasi-compact, quasi-DM, and separated.
\end{lemma}

\begin{proof}
Assume $\mathcal{X}$ is well-nigh affine. Choose an affine scheme $U$
and a surjective, flat, finite, and finitely presented morphism
$U \to \mathcal{X}$. Set $R = U \times_\mathcal{X} U$. Then we
obtain a groupoid $(U, R, s, t, c)$ in algebraic spaces and an
isomorphism $[U/R] \to \mathcal{X}$, see
Algebraic Stacks, Lemma \ref{algebraic-lemma-map-space-into-stack}
and Remark \ref{algebraic-remark-flat-fp-presentation}.
Since $s, t : R \to U$ are 
flat, finite, and finitely presented morphisms
(as base changes of $U \to \mathcal{X})$ we see that
$s, t$ are finite locally free
(Morphisms, Lemma \ref{morphisms-lemma-finite-flat}).
This implies that $R$ is affine (as finite morphisms are affine)
and hence (2) holds.

\medskip\noindent
Suppose that we have a groupoid scheme $(U, R, s, t, c)$ with $U$ and
$R$ are affine and $s, t : R \to U$ finite locally free.
Set $\mathcal{X} = [U/R]$. Then $\mathcal{X}$ is an algebraic stack
by Criteria for Representability, Theorem
\ref{criteria-theorem-flat-groupoid-gives-algebraic-stack} (strictly speaking
we don't need this here, but it can't be stressed enough that this is true).
The morphism $U \to \mathcal{X}$ is surjective, flat, and locally
of finite presentation by
Criteria for Representability, Lemma
\ref{criteria-lemma-flat-quotient-flat-presentation}.
Thus we can check whether $U \to \mathcal{X}$ is finite by
checking whether the projection $U \times_\mathcal{X} U \to U$
has this property, see Properties of Stacks, Lemma
\ref{stacks-properties-lemma-check-property-covering}.
Since $U \times_\mathcal{X} U = R$ by
Groupoids in Spaces, Lemma
\ref{spaces-groupoids-lemma-quotient-stack-2-cartesian}
we see that this is true. Thus $\mathcal{X}$ is well-nigh affine.

\medskip\noindent
Proof of the final statement. We see that $\mathcal{X}$
is quasi-compact by Properties of Stacks, Lemma
\ref{stacks-properties-lemma-quasi-compact-stack}.
We see that $\mathcal{X} = [U/R]$ is quasi-DM and separated by
Morphisms of Stacks, Lemma
\ref{stacks-morphisms-lemma-properties-diagonal-from-presentation}.
\end{proof}

\begin{lemma}
\label{lemma-affine-over-well-nigh-affine}
Let the algebraic stack $\mathcal{X}$ be well-nigh affine.
\begin{enumerate}
\item If $\mathcal{X}$ is an algebraic space, then it is affine.
\item If $\mathcal{X}' \to \mathcal{X}$ is an affine morphism
of algebraic stacks, then $\mathcal{X}'$ is well-nigh affine.
\end{enumerate}
\end{lemma}

\begin{proof}
Part (1) follows from immediately from
Limits of Spaces, Lemma \ref{spaces-limits-lemma-affine}.
However, this is overkill, since (1) also follows from
Lemma \ref{lemma-well-nigh-affine} combined with
Groupoids, Proposition
\ref{groupoids-proposition-finite-flat-equivalence}.

\medskip\noindent
To prove (2) we choose an affine scheme $U$ and a
surjective, flat, finite, and finitely presented morphism $U \to \mathcal{X}$.
Then $U' = \mathcal{X}' \times_\mathcal{X} U$ admits an affine
morphism to $U$ (Morphisms of Stacks, Lemma
\ref{stacks-morphisms-lemma-base-change-affine}).
Therefore $U'$ is an affine scheme. Of course
$U' \to \mathcal{X}'$ is surjective, flat, finite, and finitely presented
as a base change of $U \to \mathcal{X}$.
\end{proof}

\begin{lemma}
\label{lemma-well-nigh-affine-moduli-space}
Let the algebraic stack $\mathcal{X}$ be well-nigh affine. There exists
a uniform categorical moduli space
$$
f : \mathcal{X} \longrightarrow M
$$
in the category of affine schemes. Moreover
$f$ is separated, quasi-compact, and a universal homeomorphism.
\end{lemma}

\begin{proof}
Write $\mathcal{X} = [U/R]$ with $(U, R, s, t, c)$ as in
Lemma \ref{lemma-well-nigh-affine}. Let $C$ be the ring of
$R$-invariant functions on $U$, see
Groupoids, Section \ref{groupoids-section-finite-flat}.
We set $M = \Spec(C)$. The $R$-invariant morphism
$U \to M$ corresponds to a morphism $f : \mathcal{X} \to M$ by
Lemma \ref{lemma-quotient-compare}.
The characterization of morphisms into affine schemes given in
Schemes, Lemma \ref{schemes-lemma-morphism-into-affine}
immediately guarantees that $\phi : U \to M$ is a categorical
quotient in the category of affine schemes. Hence $f$ is a
categorical moduli space in the category of affine schemes
(Lemma \ref{lemma-categorical-quotient-compare}).

\medskip\noindent
Since $\mathcal{X}$ is separated by Lemma \ref{lemma-well-nigh-affine}
we find that $f$ is separated by Morphisms of Stacks, Lemma
\ref{stacks-morphisms-lemma-compose-after-separated}.

\medskip\noindent
Since $U \to \mathcal{X}$ is surjective and since $U \to M$ is quasi-compact,
we see that $f$ is quasi-compact by Morphisms of Stacks, Lemma
\ref{stacks-morphisms-lemma-surjection-from-quasi-compact}.

\medskip\noindent
By Groupoids, Lemma \ref{groupoids-lemma-integral-over-invariants}
the composition
$$
U \to \mathcal{X} \to M
$$
is an integral morphism of affine schemes. In particular, it is
universally closed
(Morphisms, Lemma \ref{morphisms-lemma-integral-universally-closed}).
Since $U \to \mathcal{X}$ is surjective, it follows that $\mathcal{X} \to M$
is universally closed (Morphisms of Stacks, Lemma
\ref{stacks-morphisms-lemma-image-proper-is-proper}).
To conclude that $\mathcal{X} \to M$ is a universal homeomorphism,
it is enough to show that it is universally bijective, i.e.,
surjective and universally injective.

\medskip\noindent
We have $|\mathcal{X}| = |U|/|R|$ by
Morphisms of Stacks, Lemma \ref{stacks-morphisms-lemma-points-presentation}.
Thus $|f|$ is surjective and even bijective
by Groupoids, Lemma \ref{groupoids-lemma-points}.

\medskip\noindent
Let $C \to C'$ be a ring map. Let $(U', R', s', t', c')$ be
the base change of $(U, R, s, t, c)$ by $M' = \Spec(C') \to M$.
Setting $\mathcal{X}' = [U'/R']$, we observe that
$M' \times_M \mathcal{X} = \mathcal{X}'$ by
Quotients of Groupoids, Lemma
\ref{groupoids-quotients-lemma-base-change-quotient-stack}.
Let $C^1$ be the ring of $R'$-invariant functions on $U'$.
Set $M^1 = \Spec(C^1)$ and consider the diagram
$$
\xymatrix{
\mathcal{X}' \ar[d]^{f'} \ar[r] & \mathcal{X} \ar[dd]^f \\
M^1 \ar[d] \\
M' \ar[r] & M
}
$$
By Groupoids, Lemma \ref{groupoids-lemma-invariants-base-change} and
Algebra, Lemma \ref{algebra-lemma-universally-bijective}
the morphism $M^1 \to M'$ is a homeomorphism.
On the other hand, the previous paragraph applied to
$(U', R', s', t', c')$ shows that $|f'|$ is bijective.
We conclude that $f$ induces a bijection on points after any
base change by an affine scheme. Thus $f$ is universally injective
by Morphisms of Stacks, Lemma
\ref{stacks-morphisms-lemma-universally-injective-local}.

\medskip\noindent
Finally, we still have to show that $f$ is a uniform moduli space
in the category of affine schemes. This follows from the discussion
above and the fact that if the
ring map $C \to C'$ is flat, then $C' \to C^1$ is an isomorphism
by Groupoids, Lemma \ref{groupoids-lemma-invariants-base-change}.
\end{proof}

\begin{lemma}
\label{lemma-well-nigh-affine-moduli-space-etale}
Let $h : \mathcal{X}' \to \mathcal{X}$ be a morphism of algebraic stacks.
Assume $\mathcal{X}'$ and $\mathcal{X}$ are well-nigh affine,
$h$ is \'etale, and $h$ induces isomorphisms on automorphism groups
(Morphisms of Stacks, Remark
\ref{stacks-morphisms-remark-identify-automorphism-groups}).
Then there exists a cartesian diagram
$$
\xymatrix{
\mathcal{X}' \ar[d] \ar[r] & \mathcal{X} \ar[d] \\
M' \ar[r] & M
}
$$
where $M' \to M$ is \'etale and
the vertical arrows are the moduli spaces constructed in
Lemma \ref{lemma-well-nigh-affine-moduli-space}.
\end{lemma}

\begin{proof}
Observe that $h$ is representable by algebraic spaces by
Morphisms of Stacks, Lemmas
\ref{stacks-morphisms-lemma-aut-iso-unramified} and
\ref{stacks-morphisms-lemma-stabilizer-preserving}.
Choose an affine scheme $U$ and a
surjective, flat, finite, and finitely presented morphism $U \to \mathcal{X}$.
Then $U' = \mathcal{X}' \times_\mathcal{X} U$ is an algebraic
space with a finite (in particular affine) morphism $U' \to \mathcal{X}'$.
By Lemma \ref{lemma-affine-over-well-nigh-affine}
we conclude that $U'$ is affine.
Setting $R = U \times_\mathcal{X} U$ and $R' = U' \times_{\mathcal{X}'} U'$
we obtain groupoids $(U, R, s, t, c)$ and $(U', R', s', t', c')$
such that $\mathcal{X} = [U/R]$ and $\mathcal{X}' = [U'/R']$,
see proof of Lemma \ref{lemma-well-nigh-affine}.
we see that the diagrams
$$
\xymatrix{
R' \ar[d]_{s'} \ar[r]_f & R \ar[d]^s \\
U' \ar[r]^f & U
}
\quad
\quad
\xymatrix{
R' \ar[d]_{t'} \ar[r]_f & R \ar[d]^t \\
U' \ar[r]^f & U
}
\quad
\quad
\xymatrix{
G' \ar[d] \ar[r]_f & G \ar[d] \\
U' \ar[r]^f & U
}
$$
are cartesian where $G$ and $G'$ are the stabilizer group schemes.
This follows for the first two by transitivity of fibre products
and for the last one this follows because it is the pullback of the
isomorphism $\mathcal{I}_{\mathcal{X}'} \to
\mathcal{X}' \times_\mathcal{X} \mathcal{I}_\mathcal{X}$
(by the already used Morphisms of Stacks, Lemma
\ref{stacks-morphisms-lemma-aut-iso-unramified}).
Recall that $M$, resp.\ $M'$ was constructed in
Lemma \ref{lemma-well-nigh-affine-moduli-space}
as the spectrum of the ring of $R$-invariant functions on $U$,
resp.\ the ring of $R'$-invariant functions on $U'$.
Thus $M' \to M$ is \'etale and $U' = M' \times_M U$
by Groupoids, Lemma \ref{groupoids-lemma-etale}.
It follows that $R' = M' \times_M U$, in other words
the groupoid $(U', R', s', t', c')$ is the base change of
$(U, R, s, t, c)$ by $M' \to M$.
This implies that the diagram in the lemma is
cartesian by
Quotients of Groupoids, Lemma
\ref{groupoids-quotients-lemma-base-change-quotient-stack}.
\end{proof}

\begin{lemma}
\label{lemma-moduli-space-finite-affine}
Let the algebraic stack $\mathcal{X}$ be well-nigh affine. The morphism
$$
f : \mathcal{X} \longrightarrow M
$$
of Lemma \ref{lemma-well-nigh-affine-moduli-space}
is a uniform categorical moduli space.
\end{lemma}

\begin{proof}
We already know that $M$ is a uniform categorical moduli space
in the category of affine schemes. By
Lemma \ref{lemma-check-uniform-categorical-quotient-on-affines}
it suffices to show that the base change
$f' : M' \times_M \mathcal{X} \to M'$
is a categorical moduli space for any flat morphism
$M' \to M$ of affine schemes.
Observe that $\mathcal{X}' = M' \times_M \mathcal{X}$ is well-nigh affine by
Lemma \ref{lemma-affine-over-well-nigh-affine}.
This after replacing $\mathcal{X}$ by $\mathcal{X}'$
and $M$ by $M'$, we reduce to proving $f$ is a categorical
moduli space.

\medskip\noindent
Let $g : \mathcal{X} \to Y$ be a morphism where $Y$ is an algebraic space.
We have to show that $g = h \circ f$ for a unique morphism $h : M \to Y$.

\medskip\noindent
Uniqueness. Suppose we have two morphisms $h_i : M \to Y$ with
$g = h_1 \circ f = h_2 \circ f$. Let $M' \subset M$ be the equalizer
of $h_1$ and $h_2$. Then $M' \to M$ is a monomorphism and
$f : \mathcal{X} \to M$ factors through $M'$. Thus $M' \to M$
is a universal homeomorphism. We conclude $M'$ is affine
(Morphisms, Lemma \ref{morphisms-lemma-universal-homeomorphism}).
But then since $f : \mathcal{X} \to M$
is a categorical moduli space in the category of affine schemes,
we see $M' = M$.

\medskip\noindent
Existence. Below we will show that for every $p \in M$ there exists
a cartesian square
$$
\xymatrix{
\mathcal{X}' \ar[r] \ar[d] & \mathcal{X} \ar[d] \\
M' \ar[r] & M
}
$$
with $M' \to M$ an \'etale morphism of affines and $p$ in the image such that
the composition $\mathcal{X}' \to \mathcal{X} \to Y$ factors through $M'$.
This means we can construct the map $h : M \to Y$ \'etale locally on $M$.
Since $Y$ is a sheaf for the \'etale topology and by the uniqueness shown
above, this is enough (small detail omitted).

\medskip\noindent
Let $y \in |Y|$ be the image of $p$.
Let $(V, v) \to (Y, y)$ be an \'etale morphism with $V$ affine.
Consider $\mathcal{X}' = V \times_Y \mathcal{X}$.
Observe that $\mathcal{X}' \to \mathcal{X}$ is separated and \'etale
as the base change of $V \to Y$. Moreover, $\mathcal{X}' \to \mathcal{X}$
induces isomorphisms on automorphism groups
(Morphisms of Stacks, Remark
\ref{stacks-morphisms-remark-identify-automorphism-groups})
as this is true for
$V \to Y$, see Morphisms of Stacks, Lemma
\ref{stacks-morphisms-lemma-base-change-stabilizer-preserving}.
Choose a presentation $\mathcal{X} = [U/R]$
as in Lemma \ref{lemma-well-nigh-affine}.
Set $U' = \mathcal{X}' \times_\mathcal{X} U = V \times_Y U$
and choose $u' \in U'$ mapping to $p$ and $v$ (possible
by Properties of Spaces, Lemma \ref{spaces-properties-lemma-points-cartesian}).
Since $U' \to U$ is separated and \'etale we see that
every finite set of points of $U'$ is contained in an affine open, see
More on Morphisms, Lemma
\ref{more-morphisms-lemma-separated-locally-quasi-finite-over-affine}.
On the other hand, the morphism $U' \to \mathcal{X}'$ is
surjective, finite, flat, and locally of finite presentation.
Setting $R' = U' \times_{\mathcal{X}'} U'$ we see
that $s', t' : R' \to U'$ are finite locally free.
By Groupoids, Lemma \ref{groupoids-lemma-find-invariant-affine}
there exists an $R'$-invariant affine open subscheme $U'' \subset U'$
containing $u'$.
Let $\mathcal{X}'' \subset \mathcal{X}'$ be
the corresponding open substack. Then $\mathcal{X}''$ is
well-nigh affine. By Lemma \ref{lemma-well-nigh-affine-moduli-space-etale}
we obtain a cartesian square
$$
\xymatrix{
\mathcal{X}'' \ar[r] \ar[d] & \mathcal{X} \ar[d] \\
M'' \ar[r] & M
}
$$
with $M'' \to M$ \'etale. Since $\mathcal{X}'' \to M''$ is
a categorical moduli space in the category of affine schemes
we obtain a morphism $M'' \to V$ such that the composition
$\mathcal{X}'' \to \mathcal{X}' \to V$ is equal to the composition
$\mathcal{X}'' \to M'' \to V$. This proves our claim and finishes
the proof.
\end{proof}

\begin{lemma}
\label{lemma-etale-separated-over-well-nigh-affine}
Let $h : \mathcal{X}' \to \mathcal{X}$ be a morphism of algebraic stacks.
Assume $\mathcal{X}$ is well-nigh affine, $h$ is \'etale, $h$ is separated,
and $h$ induces isomorphisms on automorphism groups
(Morphisms of Stacks, Remark
\ref{stacks-morphisms-remark-identify-automorphism-groups}).
Then there exists a cartesian diagram
$$
\xymatrix{
\mathcal{X}' \ar[d] \ar[r] & \mathcal{X} \ar[d] \\
M' \ar[r] & M
}
$$
where $M' \to M$ is a separated \'etale morphism of schemes and
$\mathcal{X} \to M$ is the moduli space constructed in
Lemma \ref{lemma-well-nigh-affine-moduli-space}.
\end{lemma}

\begin{proof}
Choose an affine scheme $U$ and a surjective, flat, finite, and
locally finitely presented morphism $U \to \mathcal{X}$.
Since $h$ is representable by algebraic spaces
(Morphisms of Stacks, Lemmas
\ref{stacks-morphisms-lemma-aut-iso-unramified} and
\ref{stacks-morphisms-lemma-stabilizer-preserving})
we see that $U' = \mathcal{X}' \times_\mathcal{X} U$ is
an algebraic space. Since $U' \to U$ is separated and \'etale,
we see that $U'$ is a scheme and that every finite set of points
of $U'$ is contained in an affine open, see
Morphisms of Spaces, Lemma
\ref{spaces-morphisms-lemma-locally-quasi-finite-separated-representable}
and
More on Morphisms, Lemma
\ref{more-morphisms-lemma-separated-locally-quasi-finite-over-affine}.
Setting $R' = U' \times_{\mathcal{X}'} U'$ we see
that $s', t' : R' \to U'$ are finite locally free.
By Groupoids, Lemma \ref{groupoids-lemma-find-invariant-affine}
there exists an open covering $U' = \bigcup U'_i$ by
$R'$-invariant affine open subschemes $U'_i \subset U'$.
Let $\mathcal{X}'_i \subset \mathcal{X}'$ be the corresponding
open substacks. These are well-nigh affine as $U'_i \to \mathcal{X}'_i$
is surjective, flat, finite and of finite presentation. By
Lemma \ref{lemma-well-nigh-affine-moduli-space-etale}
we obtain cartesian diagrams
$$
\xymatrix{
\mathcal{X}'_i \ar[r] \ar[d] & \mathcal{X} \ar[d] \\
M'_i \ar[r] & M
}
$$
with $M'_i \to M$ an \'etale morphism of affine schemes
and vertical arrows as in
Lemma \ref{lemma-well-nigh-affine-moduli-space}.
Observe that
$\mathcal{X}'_{ij} = \mathcal{X}'_i \cap \mathcal{X}'_j$
is an open subspace of $\mathcal{X}'_i$ and $\mathcal{X}'_j$.
Hence we get corresponding open subschemes
$V_{ij} \subset M'_i$ and $V_{ji} \subset M'_j$.
By the result of
Lemma \ref{lemma-moduli-space-finite-affine}
we see that both
$\mathcal{X}'_{ij} \to V_{ij}$ and
$\mathcal{X}'_{ji} \to V_{ji}$ are categorical moduli spaces!
Thus we get a unique isomorphism $\varphi_{ij} : V_{ij} \to V_{ji}$
such that
$$
\xymatrix{
\mathcal{X}'_i \ar[d] & &
\mathcal{X}'_i \cap \mathcal{X}'_j \ar[rr] \ar[ll] \ar[ld] \ar[rd] & &
\mathcal{X}'_j \ar[d] \\
M'_i &
V_{ij} \ar[l] \ar[rr]^{\varphi_{ij}} & &
V_{ji} \ar[r] &
M'_j
}
$$
is commutative. These isomorphisms satisfy the cocyclce condition of
Schemes, Section \ref{schemes-section-glueing-schemes} by a computation
(and another application of the previous lemma) which we omit.
Thus we can glue the affine schemes in to scheme $M'$, see
Schemes, Lemma \ref{schemes-lemma-glue}.
Let us identify the $M'_i$ with their image in $M'$.
We claim there is a morphism $\mathcal{X}' \to M'$ fitting into
cartesian diagrams
$$
\xymatrix{
\mathcal{X}'_i \ar[r] \ar[d] & \mathcal{X}' \ar[d] \\
M'_i \ar[r] & M'
}
$$
This is clear from the description of the morphisms into the glued scheme $M'$
in Schemes, Lemma \ref{schemes-lemma-glue} and the fact that to give a morphism
$\mathcal{X}' \to M'$ is the same thing as given a morphism $T \to M'$
for any morphism $T \to \mathcal{X}'$.
Similarly, there is a morphism $M' \to M$ restricting to the
given morphisms $M'_i \to M$ on $M'_i$.
The morphism $M' \to M$ is \'etale (being \'etale on the members of an
\'etale covering) and the fibre product property holds as it can
be checked on members of the (affine) open covering $M' = \bigcup M'_i$.
Finally, $M' \to M$ is separated because the composition
$U' \to \mathcal{X}' \to M'$ is surjective and universally closed
and we can apply Morphisms, Lemma
\ref{morphisms-lemma-image-universally-closed-separated}.
\end{proof}

\begin{lemma}
\label{lemma-etale-local-finite-inertia}
Let $\mathcal{X}$ be an algebraic stack. Assume
$\mathcal{I}_\mathcal{X} \to \mathcal{X}$ is finite.
Then there exist a set $I$ and for $i \in I$ a morphism of algebraic stacks
$$
g_i : \mathcal{X}_i \longrightarrow \mathcal{X}
$$
with the following properties
\begin{enumerate}
\item $|\mathcal{X}| = \bigcup |g_i|(|\mathcal{X}_i|)$,
\item $\mathcal{X}_i$ is well-nigh affine,
\item $\mathcal{I}_{\mathcal{X}_i} \to
\mathcal{X}_i \times_\mathcal{X} \mathcal{I}_\mathcal{X}$
is an isomorphism, and
\item $g_i : \mathcal{X}_i \to \mathcal{X}$ is representable
by algebraic spaces, separated, and \'etale,
\end{enumerate}
\end{lemma}

\begin{proof}
For any $x \in |\mathcal{X}|$ we can choose
$g : \mathcal{U} \to \mathcal{X}$, $\mathcal{U} = [U/R]$, and $u$ as in
Morphisms of Stacks,
Lemma \ref{stacks-morphisms-lemma-etale-local-quasi-DM-at-x}.
Then by
Morphisms of Stacks, Lemma
\ref{stacks-morphisms-lemma-stabilizer-preserving-unramified}
we see that there exists an open substack
$\mathcal{U}' \subset \mathcal{U}$ containing $u$
such that $\mathcal{I}_{\mathcal{U}'} \to
\mathcal{U}' \times_\mathcal{X} \mathcal{I}_\mathcal{X}$
is an isomorphism.
Let $U' \subset U$ be the $R$-invariant open corresponding to
the open substack $\mathcal{U}'$.
Let $u' \in U'$ be a point of $U'$ mapping to $u$.
Observe that $t(s^{-1}(\{u'\}))$ is finite as $s : R \to U$ is finite.
By Properties, Lemma \ref{properties-lemma-ample-finite-set-in-affine}
and Groupoids, Lemma \ref{groupoids-lemma-find-invariant-affine}
we can find an $R$-invariant affine open $U'' \subset U'$
containing $u'$. Let $R''$ be the restriction of $R$ to $U''$.
Then $\mathcal{U}'' = [U''/R'']$ is an open substack of
$\mathcal{U}'$ containing $u$, is well-nigh affine,
$\mathcal{I}_{\mathcal{U}''} \to
\mathcal{U}'' \times_\mathcal{X} \mathcal{I}_\mathcal{X}$
is an isomorphism, and $\mathcal{U}'' \to \mathcal{X}$
and is representable by algebraic spaces and \'etale.
Finally, $\mathcal{U}'' \to \mathcal{X}$ is separated as
$\mathcal{U}''$ is separated (Lemma \ref{lemma-well-nigh-affine})
the diagonal of $\mathcal{X}$ is separated
(Morphisms of Stacks, Lemma \ref{stacks-morphisms-lemma-diagonal-diagonal})
and separatedness follows from Morphisms of Stacks, Lemma
\ref{stacks-morphisms-lemma-compose-after-separated}.
Since the point $x \in |\mathcal{X}|$ is arbitrary the proof is complete.
\end{proof}

\begin{proof}
We choose a set $I$\footnote{The reader who is still keeping
track of set theoretic issues should make sure $I$ is not too large.}
and for $i \in I$ a morphism of algebraic stacks
$g_i : \mathcal{X}_i \to \mathcal{X}$ as in
Lemma \ref{lemma-etale-local-finite-inertia}; we will use all
of the properties listed in this lemma without further mention.
Let
$$
f_i : \mathcal{X}_i \to M_i
$$
be as in Lemma \ref{lemma-well-nigh-affine-moduli-space}.
Consider the stacks
$$
\mathcal{X}_{ij} = \mathcal{X}_i \times_{g_i, \mathcal{X}, g_j} \mathcal{X}_j
$$
for $i, j \in I$. The projections $\mathcal{X}_{ij} \to \mathcal{X}_i$
and $\mathcal{X}_{ij} \to \mathcal{X}_j$ are separated
by Morphisms of Stacks, Lemma
\ref{stacks-morphisms-lemma-base-change-separated},
\'etale by Morphisms of Stacks, Lemma
\ref{stacks-morphisms-lemma-base-change-etale},
and induce isomorphisms on automorphism groups
(as in Morphisms of Stacks, Remark
\ref{stacks-morphisms-remark-identify-automorphism-groups}) by
Morphisms of Stacks, Lemma
\ref{stacks-morphisms-lemma-base-change-stabilizer-preserving}.
Thus we may apply Lemma \ref{lemma-etale-separated-over-well-nigh-affine}
to find a commutative diagram
$$
\xymatrix{
\mathcal{X}_i \ar[d]_{f_i} &
\mathcal{X}_{ij} \ar[d]_{f_{ij}} \ar[l] \ar[r] &
\mathcal{X}_j \ar[d]_{f_j} \\
M_i &
M_{ij} \ar[l] \ar[r] &
M_j
}
$$
with cartesian squares where $M_{ij} \to M_i$ and $M_{ij} \to M_j$
are separated \'etale morphisms of schemes; here we also use that $f_i$
is a uniform categorical quotient by
Lemma \ref{lemma-moduli-space-finite-affine}.
Claim:
$$
\coprod M_{ij} \longrightarrow \coprod M_i \times \coprod M_i
$$
is an \'etale equivalence relation.

\medskip\noindent
Proof of the claim. Set $R = \coprod M_{ij}$ and $U = \coprod M_i$.
We have already seen that $t : R \to U$ and $s : R \to U$ are \'etale.
Let us construct a morphism $c : R \times_{s, U, t} R \to R$
compatible with $\text{pr}_{13} : U \times U \times U \to U \times U$.
Namely, for $i, j, k \in I$ we consider
$$
\mathcal{X}_{ijk} =
\mathcal{X}_i \times_{g_i, \mathcal{X}, g_j} \mathcal{X}_j
\times_{g_j, \mathcal{X}, g_k} \mathcal{X}_k =
\mathcal{X}_{ij} \times_{\mathcal{X}_j} \mathcal{X}_{jk}
$$
Arguing exactly as in the previous paragraph,
we find that $M_{ijk} = M_{ij} \times_{M_j} M_{jk}$
is a categorical moduli space for $\mathcal{X}_{ijk}$.
In particular, there is a canonical morphism
$M_{ijk} = M_{ij} \times_{M_j} M_{jk} \to M_{ik}$
coming from the projection $\mathcal{X}_{ijk} \to \mathcal{X}_{ik}$.
Putting these morphisms together we obtain the morphism $c$.
In a similar fashion we construct a morphism $e : U \to R$
compatible with $\Delta : U \to U \times U$ and
$i : R \to R$ compatible with the flip $U \times U \to U \times U$.
Let $k$ be an algebraically closed field. Then
$$
\Mor(\Spec(k), \mathcal{X}_i) \to \Mor(\Spec(k), M_i) = M_i(k)
$$
is bijective on isomorphism classes and the same remains true after any
base change by a morphism $M' \to M$. This follows from our choice
of $f_i$ and Morphisms of Stacks, Lemmas
\ref{stacks-morphisms-lemma-universally-injective} and
\ref{stacks-morphisms-lemma-universally-injective-point}.
By construction of $2$-fibred products the diagram
$$
\xymatrix{
\Mor(\Spec(k), \mathcal{X}_{ij}) \ar[d] \ar[r] &
\Mor(\Spec(k), \mathcal{X}_j) \ar[d] \\
\Mor(\Spec(k), \mathcal{X}_i) \ar[r] &
\Mor(\Spec(k), \mathcal{X})
}
$$
is a fibre product of categories. By our choice of $g_i$ the
functors in this diagram induce bijections on automorphism groups.
It follows that this diagram induces a fibre product diagram
on sets of isomorphism classes! Thus we see that
$$
R(k) = U(k) \times_{|\Mor(\Spec(k), \mathcal{X})|} U(k)
$$
where $|\Mor(\Spec(k), \mathcal{X})|$ denotes the set
of isomorphism classes.
In particular, for any algebraically closed field $k$
the map on $k$-valued point is an equivalence relation.
We conclude the claim holds by
Groupoids, Lemma \ref{groupoids-lemma-etale-equivalence-relation}.

\medskip\noindent
Let $M = U/R$ be the algebraic space which is the quotient of the above
\'etale equivalence relation, see
Spaces, Theorem \ref{spaces-theorem-presentation}.
There is a canonical morphism $f : \mathcal{X} \to M$
fitting into commutative diagrams
\begin{equation}
\label{equation-fundamental-diagram}
\xymatrix{
\mathcal{X}_i \ar[r]_{g_i} \ar[d]_{f_i} & \mathcal{X} \ar[d]^f \\
M_i \ar[r] & M
}
\end{equation}
Namely, such a morphism $f$ is given by a functor
$$
f : \Mor(T, \mathcal{X}) \longrightarrow \Mor(T, M)
$$
for any scheme $T$ compatible with base change. Let $a : T \to \mathcal{X}$
be an object of the left hand side. We obtain an \'etale covering
$\{T_i \to T\}$ with $T_i = \mathcal{X}_i \times_\mathcal{X} T$
and morphisms $a_i : T_i \to \mathcal{X}_i$. Then we get
$b_i = f_i \circ a_i : T_i \to M_i$. Since
$T_i \times_T T_j = \mathcal{X}_{ij} \times_\mathcal{X} T$
we moreover get a morphism $a_{ij} : T_i \times_T T_j \to \mathcal{X}_{ij}$.
Setting $b_{ij} = f_{ij} \circ a_{ij}$ we find that
$b_i \times b_j$ factors through the monomorphism
$M_{ij} \to M_i \times M_j$. Hence the morphisms
$$
T_i \xrightarrow{b_i} M_i \to M
$$
agree on $T_i \times_T T_j$. As $M$ is a sheaf for the \'etale
topology, we see that these morphisms glue to a unique
morphism $b = f(a) : T \to M$. We omit the verification that
this construction is compatible with base change and we omit
the verification that the diagrams
(\ref{equation-fundamental-diagram}) commute.

\medskip\noindent
Claim: the diagrams (\ref{equation-fundamental-diagram}) are cartesian.
To see this we study the induced morphism
$$
h_i : \mathcal{X}_i \longrightarrow M_i \times_M \mathcal{X}
$$
This is a morphism of stacks \'etale over $\mathcal{X}$
and hence $h_i$ is \'etale (Morphisms of Stacks, Lemma
\ref{stacks-morphisms-lemma-etale-permanence}).
Since $g_i$ is separated, we see $h_i$ is separated
(use Morphisms of Stacks, Lemma
\ref{stacks-morphisms-lemma-compose-after-separated} and the fact
seen above that the diagonal of $\mathcal{X}$ is separated).
The morphism $h_i$ induces isomorphisms on automorphism groups
(Morphisms of Stacks, Remark
\ref{stacks-morphisms-remark-identify-automorphism-groups})
as this is true for $g_i$. For an algebraically closed field $k$
the diagram
$$
\xymatrix{
\Mor(\Spec(k), M_i \times_M \mathcal{X}) \ar[r] \ar[d] &
\Mor(\Spec(k), \mathcal{X}) \ar[d] \\
M_i(k) \ar[r] &
M(k)
}
$$
is a catesian diagram of categories and the top arrow
induces bijections on automorphism groups.
On the other hand, we have
$$
M(k) = U(k)/R(k) = U(k)/
U(k) \times_{|\Mor(\Spec(k), \mathcal{X})|} U(k) =
|\Mor(\Spec(k), \mathcal{X})|
$$
by what we said above. Thus the right vertical arrow in the
cartesian diagram above is a bijection on isomorphism classes.
We conclude that
$|\Mor(\Spec(k), M_i \times_M \mathcal{X})| \to M_i(k)$ is bijective.
Review: $h_i$ is a separated, \'etale, induces isomorphisms on
automorphism groups (as in Morphisms of Stacks, Remark
\ref{stacks-morphisms-remark-identify-automorphism-groups}), and
induces an equivalence on fibre categories over algebraically closed fields.
Hence it is an isomorphism by Morphisms of Stacks, Lemma
\ref{stacks-morphisms-lemma-etale-iso}.

\medskip\noindent
From the claim we get in particular the following:
we have a surjective \'etale morphism $U \to M$
such that the base change of $f$ is separated, quasi-compact,
and a universal homeomorphism. It follows that $f$ is separated,
quasi-compact, and a universal homeomorphism.
See Morphisms of Stacks, Lemma
\ref{stacks-morphisms-lemma-check-separated-covering},
\ref{stacks-morphisms-lemma-check-quasi-compact-covering}, and
\ref{stacks-morphisms-lemma-check-universal-homeomorphism-covering}

\medskip\noindent
To finish the proof we have to show that $f : \mathcal{X} \to M$
is a uniform categorical moduli space.
To prove this it suffices to show that given a flat morphism
$M' \to M$ of algebraic spaces, the base change
$$
M' \times_M \mathcal{X} \longrightarrow M'
$$
is a categorical moduli space. Thus we consider a morphism
$$
\theta : M' \times_M \mathcal{X} \longrightarrow E
$$
where $E$ is an algebraic space. For each $i$ we know that
$f_i$ is a uniform categorical moduli space. Hence we obtain
$$
\xymatrix{
M' \times_M \mathcal{X}_i \ar[d] \ar[r] &
M' \times_M \mathcal{X} \ar[d]^\theta \\
M' \times_M M_i \ar[r]^{\psi_i} &
E
}
$$
Since $\{M' \times_M M_i \to M'\}$ is an \'etale covering,
to obtain the desired morphism $\psi : M' \to E$ it suffices
to show that $\psi_i$ and $\psi_j$ agree over
$M' \times_M M_i \times_M M_j = M' \times_M M_{ij}$.
This follows easily from the fact that
$f_{ij} : \mathcal{X}_{ij} =
\mathcal{X}_i \times_\mathcal{X} \mathcal{X}_j \to M_{ij}$ is a uniform
categorical quotient; details omitted.
Then finally one shows that $\psi$ fits into the commutative diagram
$$
\xymatrix{
M' \times_M \mathcal{X} \ar[d] \ar[rd]^\theta \\
M' \ar[r]^\psi &
E
}
$$
because ``$\{M' \times_M \mathcal{X}_i \to M' \times_M \mathcal{X}\}$
is an \'etale covering'' and the morphisms $\psi_i$ fit into the
corresponding commutative diagrams by construction.
This finishes the proof of the Keel-Mori theorem.
\end{proof}

\bibliographystyle{plain}
\bibliography{my}
\end{document}
